% !TeX program = pdflatex
\documentclass[10pt,a4paper]{article}

\usepackage[utf8]{inputenc}
\usepackage[T1]{fontenc}
\usepackage{geometry}
\usepackage[hidelinks]{hyperref}
\usepackage{xcolor}

\geometry{
  top=1.7cm,
  bottom=1.8cm,
  left=1.8cm,
  right=1.8cm
}

\pagestyle{empty}
\frenchspacing
\setlength{\parindent}{0pt}
\setlength{\parskip}{0.45em}

\definecolor{mydarkgreen}{RGB}{0,0,0}

\begin{document}

\begin{minipage}[t]{0.48\textwidth}
\textbf{\large Javier Andres TARAZONA JIMENEZ}\\
Étudiant ingénieur IA / Data Science\\
Massy, Île-de-France, France\\
\href{mailto:javitar06@gmail.com}{javitar06@gmail.com}\\
+33 07 46 36 97 75\\
\href{https://www.linkedin.com/in/javier-andres-tarazona-jimenez-84b489222/}{LinkedIn}
\end{minipage}
\hfill
\begin{minipage}[t]{0.42\textwidth}
\raggedleft
Crédit Agricole Assurances\\
Direction des Assurances à l'International\\
Pôle Data \& Modèles\\
Paris, France
\end{minipage}

\vspace{1.1cm}

\raggedleft
Massy, le 19 juin 2026

\vspace{0.7cm}

\raggedright
\textbf{Objet : Candidature à l'alternance Data Scientist H/F}

\vspace{0.5cm}

{\small

Madame, Monsieur,

Je suis étudiant colombien en double diplôme entre l’Universidad Nacional de Colombia et ENSTA Paris, au sein de l’Institut Polytechnique de Paris, avec une spécialisation en intelligence artificielle et data science. Mon parcours en informatique s’est construit autour d’une combinaison de motivations personnelles et intellectuelles : le besoin de comprendre des systèmes complexes, le goût de transformer des idées abstraites en solutions concrètes, la recherche d’autonomie technique, ainsi qu’une forte exigence académique et compétitive.

Au fil de mon parcours, l’intelligence artificielle est apparue comme une évolution naturelle : passer de la programmation de systèmes à la conception de systèmes capables de percevoir, d’apprendre et d’aider à la prise de décision. Ce domaine m’attire particulièrement parce qu’il oblige à travailler sur des problèmes ouverts, souvent ambigus, qui demandent à la fois rigueur scientifique, pragmatisme et capacité d’adaptation. La vision par ordinateur représente pour moi un lien très concret entre mathématiques, code et perception du monde réel, tandis que l’IA, plus largement, constitue une technologie transversale avec un impact visible sur la société et l’économie.

C’est dans cette perspective que Crédit Agricole S.A. et Crédit Agricole Assurances représentent pour moi un environnement particulièrement attractif. Au-delà de la reconnaissance de Crédit Agricole Assurances comme acteur majeur du secteur financier, votre position de premier bancassureur européen et de premier assureur en France par chiffre d’affaires témoigne d’un impact réel à grande échelle. Votre stratégie ACT 2028, orientée vers le leadership européen, les transitions et les nouvelles technologies, renforce également l’intérêt de votre groupe pour un profil souhaitant développer des solutions d’IA utiles, fiables et intégrées à des enjeux industriels concrets.

Ce qui m’intéresse particulièrement dans le secteur de l’assurance est son impact humain direct. Les sujets liés à la protection, à l’épargne, à la retraite, à la santé, au logement, aux accidents ou à la stabilité familiale montrent que les décisions techniques ne sont pas uniquement abstraites : elles peuvent contribuer à améliorer des services qui accompagnent des moments importants de la vie des clients. Crédit Agricole Assurances offre ainsi un équilibre rare entre ingénierie, données, métier, responsabilité et impact.

Votre alternance représente pour moi un pont idéal entre l’IA technique et l’application industrielle réelle. Mon expérience actuelle est centrée sur la vision par ordinateur, la recherche en vision 3D, le NLP, l’IA générative et le développement logiciel. Dans le domaine de l’assurance, l’IA ne reste pas au stade de démonstration : elle peut être intégrée à des processus métier concrets, comme l’analyse documentaire, l’automatisation de reportings, l’extraction d’informations, la classification de documents, la détection d’anomalies ou encore l’amélioration de modèles de décision. Cette immersion dans un environnement financier et assurantiel me permettrait de renforcer un profil hybride associant IA, software engineering, data science, risque et compréhension métier.

Je pense pouvoir apporter à Crédit Agricole Assurances une combinaison cohérente de compétences techniques, d’autonomie et de capacité d’apprentissage. Lors de mon stage de recherche à ENSTA Paris – U2IS, je travaille sur le pré-entraînement auto-supervisé avec PyTorch et Vision Transformers, ainsi que sur la génération de jeux de données contrôlés avec Python et Blender. Chez APEDYS91, j’ai développé une application NLP intégrant des règles linguistiques, des LLMs locaux, de la reconnaissance d’entités et des contrôles de cohérence. Chez Engin A.I, j’ai contribué à des pipelines de vision par ordinateur et d’analyse de données, tout en améliorant la structure logicielle de bibliothèques Python dans une logique modulaire et maintenable.

Ces expériences correspondent directement aux besoins de votre pôle Data \& Modèles : Machine Learning, modèles supervisés et non supervisés, NLP, Computer Vision, IA générative, traitement documentaire, automatisation, Git, CI/CD et travail dans un contexte international. Mon profil international, ma capacité à communiquer en français, anglais et espagnol, ainsi que mon intérêt pour les environnements multiculturels constituent également un atout pour une direction amenée à interagir avec plusieurs filiales du groupe à l’international.

Je serais ravi de pouvoir mettre mes compétences en intelligence artificielle, data science et ingénierie logicielle au service des projets de Crédit Agricole Assurances, tout en développant une compréhension plus approfondie des enjeux de l’assurance, de la bancassurance et de la transformation technologique du secteur financier.

Je vous remercie pour l’attention portée à ma candidature et me tiens à votre disposition pour un entretien.


Veuillez agréer, Madame, Monsieur, l'expression de mes salutations distinguées.

\vspace{0.8cm}

\textbf{Javier Andrés Tarazona Jimenez}

}

\end{document}
